\section{Conclusion}
\label{sec:conclusion}

We began by noting that JTT/LfF succeed empirically but their temporal ordering assumption remained unverified---spurious features hypothesized to be learned earlier, but never measured at the gradient level. Our gradient-level measurement validates this foundation: spurious features converge 4 epochs earlier than core features on CMNIST ($E_s=13$ vs $E_c=17$, $\Delta=4$ epochs), with early convolutional layers showing significantly higher neuron-spurious correlation than late layers ($p=0.0028$, Cohen's $d=0.25$). This mechanistic evidence explains why reweighting late-learned examples (JTT) succeeds---it implicitly corrects for temporal misalignment between spurious and core feature convergence driven by gradient descent's implicit bias toward simpler features.

However, when attempting to exploit this temporal signal for gradient-aware interventions, cross-dataset transfer failed catastrophically (39\% worst-group accuracy vs 86\% JTT target). Neuron-spurious correlations $\rho_j$ are dataset-specific and spurious-feature-type dependent---CMNIST color spurious correlation produces different neuron activation patterns than Waterbirds background correlation. This negative result identifies a fundamental constraint: gradient-aware methods require per-dataset calibration, limiting deployment scalability compared to example-level reweighting.

Temporal ordering is measurable, mechanistically grounded in architectural feature hierarchy, and explains why reweighting works. Future debiasing methods can exploit this signal, but must account for dataset-specific feature representations rather than assuming universal neuron correlations. Beyond multi-dataset validation (Waterbirds, CelebA, NICO++) to test generalization beyond color-based spurious correlations, temporal gaps open possibilities for adaptive reweighting schedules (monitor gradient convergence in real-time, adjust example weights dynamically as gaps emerge) and architectural interventions (enforce late-layer learning delays via layer-wise freezing schedules to prevent early spurious dominance).

The question is not whether temporal ordering exists---our gradient-level measurement settles that---but how to design systems that inherently resist it. CNNs amplify temporal gaps via hierarchical processing; do Vision Transformers with global attention reduce gaps? Do adaptive optimizers (Adam, AdamW) equalize convergence speeds between spurious and core features? Can we predict temporal gap magnitude from dataset spurious correlation strength ($\rho_{\text{data}}=0.75$ for CMNIST) and feature complexity metrics? These questions define the research frontier---moving from validation to principled intervention design grounded in optimization dynamics.
